\documentclass[english, parskip=full]{scrartcl}
\usepackage[utf8]{inputenc}  % Kodierung der Datei
\usepackage[english]{babel}  % Sprache auf Deutsch 
\usepackage{color}
\usepackage{graphicx, gensymb}
\usepackage{amsfonts, cancel, amssymb}
\usepackage{amsmath, amsthm, array, esint}
\usepackage{booktabs, multirow}  % schönere Tabellen
\usepackage{caption}  % Anpassung der Captions
\usepackage{scrlayer-scrpage}    % Manipulation des Seitenstils
%\pagestyle{scrheadings}  % Stil für die Seite setzen
%\automark{section}  % Abschnittsnamen als Seitenbeschriftung 
%\setheadsepline{.5pt}    % Kopzeile durch Linie abtrennen
\usepackage[hidelinks]{hyperref}  % Links und weitere PDF-Features
\usepackage{typearea, tikz, pgffor, verbatim}
\usetikzlibrary{calc, quotes}
\usepackage[cache=false, outputdir=build]{minted}
\usepackage{placeins} % provides \FloatBarrier

\definecolor{bg}{rgb}{0.9,0.9,0.9}
\newcommand\numberthis{\addtocounter{equation}{1}\tag{\theequation}}
\newcommand{\bra}{}
\newcommand{\kom}{}
\newcommand{\brak}{}
\newcommand{\braket}{}
\newcommand{\intx}{}
\newcommand{\intr}{}
\newcommand{\kla}{}
\newcommand{\klam}{}
\newcommand{\klammer}{}
\newcommand{\oper}{}
\newcommand{\tecxt}{}
\newcommand{\Bra}{}
\newcommand{\Ket}{}
\renewcommand{\i}{\text{i}}
\newcommand{\e}{\text{e}}
\newcommand*{\Fourier}{\textsc{Fourier}\ }
\newcommand*{\F}{\mathcal{F}}
\newcommand*{\sinc}{\text{sinc\,}}
\newcommand{\conv}{\mathop{\scalebox{3.5}{\raisebox{-0.38ex}{\makebox[0.5\width][c]{$\ast$}}}}\limits}

\newcommand*{\titel}{06 Newton-Verfahren}
\newcommand*{\autor}{Joachim Schwardt und Julian Fleck}

\hypersetup{pdfauthor={\autor}, pdftitle={\titel}}

%\titlehead{titlehead}
\subject{Numerik I - Aufgaben}
\title{\titel}
%\author{\autor}
\author{\begin{tabular}{l|l}
Joachim Schwardt & 4768711 \\ 
Julian Fleck & 4759587\\ 
\end{tabular}}
%\author{\begin{tabular}{|l|l|}
%\hline
%Joachim Schwardt & 4768711 \\ \hline
%Julian Fleck & 4759587\\ \hline
%\end{tabular}}
\date{\today}

\begin{document}

%\begin{titlepage}
\maketitle  % Titel setzen
%\tableofcontents  % Inhaltsverzeichnis setzen
%\end{titlepage}

\section*{Aufgabe 1}
Sei $\phi : \mathbb{R} \rightarrow \mathbb{R}$ eine Funktion gegeben durch
\begin{align}
\phi(x) &= |x| + \frac{1}{1 + |x|}.
\end{align}
\subsection*{a)}
Sei $x\neq y\in\mathbb{R}$. 
Dann gilt
\begin{align}
|\phi(x) - \phi(y)| &= \left\vert |x| - |y| + \frac{1}{1 + |x|} - \frac{1}{1 + |y|} \right\vert = \\
&=  \left\vert |x| - |y| + \frac{|y| - |x| }{(1+|x|)(1+|y|)} \right\vert = \\
&= \big||x| - |y| \big| \left\vert 1 - \frac{1}{(1+|x|)(1+|y|)} \right\vert \le \\
&\le |x-y| \left\vert \frac{|x| + |y| + |xy|}{1+|x| +  |y| + |xy|} \right\vert < |x-y|,
\end{align}
also ist $\phi$ eine streng kontrahierende Abbildung.
\subsection*{b)}
Angenommen es gibt ein $c\in (0,1)$, sodass für alle  $x,y\in\mathbb{R}$ die Ungleichung
\begin{align}
|\phi(x) - \phi(y)| &\le c|x-y|
\end{align}
erfüllt ist.
Sei nun $y=0$.
Dann gilt
\begin{align}
|\phi(x) - \phi(0)| &= \left\vert |x| + \frac{1}{1 + |x|} - 1 \right\vert = \\
&= |x|\left\vert 1 - \frac{1}{1 + |x|} \right\vert = : c(x) |x - 0|.
\end{align}
Die Ungleichung ist verletzt, falls $c(x) > c$.
Daraus erhalten wir eine Bedingung für $x$, durch Umstellen erhält man 
\begin{align}
c &\overset{!}{<} c(x) = \frac{|x|}{1 + |x|}\\
\Rightarrow\ \ \ |x| &> (1+|x|)c\\
\Rightarrow\ \ \ |x| &> \frac{1}{1 - c}.
\end{align}
Für $c<1$ gibt es also immer ein $x$, sodass die Ungleichung verletzt ist.
\subsection*{c)}
Wir suchen einen Fixpunkt von $\phi$, also ein $x^\ast\in\mathbb{R}$ mit $x^\ast = \phi(x^\ast)$.
Da $\phi(x) > 0$ für alle $x\in\mathbb{R}$, muss auch $x^\ast > 0$ gelten.
Allerdings ist
\begin{align}
\phi(x^\ast) &= x^\ast + \frac{1}{1 + x^\ast} > x^\ast.
\end{align}
Es gibt also keinen Fixpunkt.
Für den Banach'schen Fixpunktsatz sind tatsächlich nicht alle Voraussetzungen erfüllt.
Zwar ist die Funktion $\phi$ Lipschitz-stetig, allerdings ist $L=1$.
Der Satz garantiert den Fixpunkt nur für $L<1$.
%TODO: Warum geht hier Banach nicht? 
%Ist die Konvergenz nur für $c<1$ garantiert?

\newpage
\section*{Aufgabe 2}
Gegeben sei $a>0$ und die Iteration
\begin{align}
x_{k+1} &:= \frac{x_k}{2} + \frac{a}{2x_k},
\end{align}
wobei wir als Startwert $x_0 > \sqrt{a}$ annehmen.
\subsection*{a)}
Sei $k=0$.
Dann gilt
\begin{align}
x_{k+1} - \sqrt{a} &= x_1 - \sqrt{a} = \frac{x_0}{2} + \frac{a}{2x_0} - \sqrt{a} = \frac{x_0^2 + a - 2\sqrt{a} x_0}{2x_0} = \frac{(x_0 - \sqrt{a})^2}{2x_0}.
\end{align}
Daraus folgt offensichtlich $x_1 - \sqrt{a} > 0$ und
\begin{align}
x_{1} - \sqrt{a} &=  \frac{1}{2} (x_0 - \sqrt{a})  \frac{(x_0 - \sqrt{a})}{x_0} = \frac{1}{2} (x_0 - \sqrt{a}) \kla\left( {1 - \frac{\sqrt{a}}{x_0}} \right) \le \frac{1}{2} (x_0 - \sqrt{a}) .
\end{align}
Der Induktionsanfang ist damit gezeigt. 
Sei also nun 
\begin{align}
0 &\le x_k - \sqrt{a} \le \frac{1}{2} (x_{k-1} - \sqrt{a}) .
\end{align}
Insbesondere ist also $x_k \ge \sqrt{a}$.
Dann gilt
\begin{align}
x_{k+1} - \sqrt{a} &= \frac{x_{k}}{2} + \frac{a}{2x_{k}} - \sqrt{a} = \frac{(x_{k} - \sqrt{a})^2}{2x_{k}}.
\end{align}
Wie zuvor folgt daraus $x_{k+1} -\sqrt{a} \ge 0$ und
\begin{align}
x_{k+1} - \sqrt{a} &= \frac{1}{2} (x_{k} - \sqrt{a})  \kla\left( {1 - \frac{\sqrt{a}}{x_{k}}} \right) \kom\overset{\substack{{x_k \ge \sqrt{a}} \\ |}}{\le} \frac{1}{2} (x_k - \sqrt{a}). 
\end{align}
Zusammengefasst gilt also für alle $k\in\mathbb{N}_0$, dass
\begin{align}
0 &\le x_{k+1} - \sqrt{a} \le \frac{1}{2}(x_k - \sqrt{a}).
\end{align}
Sei $y_k := x_k - \sqrt{a}$.
Dann gilt $y_k \ge 0$ und 
\begin{align}
y_{k+1} &\le \frac{1}{2}y_k \le \dots \le \frac{1}{2^{k+1}}y_0 \kom\overset{\substack{{k \, \rightarrow\, \infty} \\ |}}{\longrightarrow} 0.
\end{align}
Also gilt $x_k \rightarrow \sqrt{a}$ für $k \rightarrow\infty$.

\subsection*{b)}
Wir hatten bereits gezeigt, dass
\begin{align}
y_{k+1} &= x_{k+1} - \sqrt{a} = \frac{(x_k - \sqrt{a})^2}{2x_k} = \frac{y_k^2}{2(y_k + \sqrt{a})}.
\end{align}
Aus $y_k\ge0$ folgt dann direkt
\begin{align}
y_{k+1} &\le \frac{y_k^2}{2\sqrt{a}},
\end{align}
die Folge konvergiert also quadratisch.

\subsection*{c)}
Sei $f(x) = x^2 - a$.
Dann lautet die Newton-Vorschrift für das iterative Bestimmen einer Nullstelle
\begin{align}
x_{k+1} &= x_k - \frac{f(x_k)}{f'(x_k)} = x_k - \frac{x_k^2 - a}{2x_k} = \frac{x_k}{2} - \frac{a}{2x_k},
\end{align}
was mit der obigen Vorschrift übereinstimmt.

\newpage
\section*{Aufgabe 3}
Gegeben sei die Funktion $f:\mathbb{R} \rightarrow\mathbb{R}$ mit $f(x) = x^3 - 3x^2 + 1$.

\subsection*{a)}
In der Tabelle \ref{table:f of x} sind einige Funktionswerte aufgeführt.
Jeder Vorzeichenwechsel bedeutet dabei, dass es in dem entsprechenden Intervall mindestens eine Nullstelle gibt.
Es gibt also mindestens eine für $x\in(-1,0)$, mindestens eine für $x\in(0,1)$ und mindestens eine für $x\in(2,3)$.
Da $f(x)$ ein Polynom dritten Grades ist, gibt es genau drei Nullstellen (bzw. im reellen unter Umständen weniger).
Daher muss es in jedem der drei Intervalle genau eine Nullstelle geben.
\begin{table}[!ht]
\centering
\begin{tabular}{c||c|c|c|c|c}
$x$ & $-1$ & $0$ & $1$ & $2$ & $3$ \\ \hline
$f(x)$ & $-3$ & $1$ & $-1$ & $-3$ & $1$\\
\end{tabular}
\caption{Einige Funktionswerte von $f(x)$.}
\label{table:f of x}
\end{table}
Diese Nullstellen seien im Folgen mit $x_j^\ast$ bezeichnet, wobei $j=1,2,3$ das entsprechende Interval angibt.


\subsection*{b)}
%Die Suche nach den Nullstellen von $f(x)$ ist äquivalent zur Suche nach den Fixpunkten von $g(x) = f(x) + x$.
In der Vorlesung hatten wir bewiesen, dass das Newton-Verfahren quadratisch konvergiert, wenn der Startwert nahe genug an der Nullstelle liegt.
Voraussetzung hierfür war, dass $f'(x)$ in der jeweiligen Umgebung einer Nullstelle Lipschitz-stetig und invertierbar ist.
Hier ist $f'(x) = 3x^2 - 6x$, wobei $f'(x) = 0$ für $x\in\{0,2\}$.
Insbesondere ist die Lipschitz-Konstante gegeben durch
\begin{align}
L_{f'} &= \max_{x\in [-1,3]} |f''(x)| =  \max_{x\in [-1,3]} 6|x-1| = 12,
\end{align}
also ist $f'(x)$ in allen Umgebungen der Nullstellen Lipschitz-stetig.
Die Nullstellen von $f'$ liegen jeweils auf Randpunkten der Intervalle, in denen wir Nullstellen von $f$ suchen.
Wir können also für alle $j=1,2,3$ Zahlen $r_j>0$ finden, sodass die Newton-Iteration für einen beliebigen Startwert $x_0\in [x_j^\ast - r_j, x_j^\ast + r_j]$ (quadratisch) konvergiert.

%Hier garantiert der Banach'sche Fixpunktsatz deren Existenz, wenn die Lipschitz-Konstante von $g$, also $L = \max_{x\in I} g'(x)$, auf dem jeweiligen Intervall streng kleiner 1 ist.
%In der Vorlesung hatten wir bewiesen, dass eine Fixpunktiteration konvergiert, falls der Startwert $x_0$ nahe genug an am Fixpunkt $x^\ast$ liegt.
%Genauer muss $|x_0 - x^\ast| < \rho := \frac{1}{\max_{x\in U} g''(x)}$ gelten, wobei $U=(x^\ast - r, x^\ast + r)$.
%Es ist $g''(x) = 6x - 6$, und somit 

\subsection*{c)}
Die Newton-Vorschrift lautet hier
\begin{align}
x_{k+1} &= x_k - \frac{f(x_k)}{f'(x_k)} = x_k - \frac{x_k^3 - 3x_k^2 + 1}{3x_k^2 - 6x_k} = \frac{2x_k^3 - 3x_k^2 - 1}{3x_k(x_k - 2)}.
\end{align}
Für $x_0=1$ ist dann $x_1 = \frac{2}{3}$.



%\begin{thebibliography}{9}
%\bibitem{example} description, \url{https://www.exmapleurl}, day.\,month~2021.
%\end{thebibliography}

\end{document}
